% ─────────────────────────────────────────────────────────────────
\section{Conclusion}
\label{sec:conclusion}
% ─────────────────────────────────────────────────────────────────

We presented a comprehensive study of protein family classification on the
SHREC\,25 benchmark, spanning the full arc from VTK mesh preprocessing to
hierarchical graph neural network training, detailed post-hoc analysis, and
comparison against the RIConv2-based method of Barisin et al.~\cite{riconv2}.

\paragraph{Pipeline and architecture.}
A two-stage pipeline converts raw VTK surface meshes into compact SSL graphs
(1\,000 nodes, 8D features) and then into geodesic Voronoi patch assignments
(64 patches).
The HierGIN model stacks node-level GIN message passing, attention-weighted
Voronoi patch pooling, patch-level GIN, and global max readout.
The two-level pooling hierarchy is the principal driver of the $+$15.8\,pp
improvement over flat GNN baselines.

\paragraph{Within-HierGIN: potential is the dominant signal.}
Ablation experiments confirm that removing electrostatic potential degrades
macro-F1 by 11.8\,pp --- the largest single-feature impact in our ablation.
NormalPotential adds a secondary gain ($+$1.4\,pp globally, critical for
class\,61).
Training HierGIN on a 14D feature set (8D\,+\,6 globally-normalised geometric
features) achieves our best result: accuracy\,0.865, balanced accuracy\,0.812,
macro-F1\,\textbf{0.808}.

\paragraph{Architecture gap: geometry is the primary signal.}
Comparison with RIConv2~\cite{riconv2} --- an independently published SHREC\,25
entry --- reveals that the primary bottleneck is architecture, not features.
RIConv2 RUN1 (xyz\,+\,normals only, \emph{no chemical features}) achieves
macro-F1\,=\,0.838 and balanced accuracy\,=\,0.813, outperforming all
HierGIN configurations including those with electrostatic features.
Adding Potential to RIConv2 (RUN2) gives only $+$1.6\,pp additional gain
(0.838$\to$0.854).
The critical difference is RIConv2's multi-scale ball-query grouping with
rotation-invariant local reference frames, which builds hierarchical local
surface descriptors invisible to a kNN-6 graph.

\paragraph{The confusion pairs are architecture-limited.}
We initially identified (61,\,90) as near-inseparable based on low marginal
KS/SW statistics for single-scale edge distributions (SW(RI)\,=\,0.004).
RIConv2 RUN1 achieves F1\,=\,0.887 on class\,61 using only geometry,
refuting this claim.
The confusion was not data-intrinsic but a consequence of HierGIN's
coarse single-scale graph representation.
The separability analysis correctly characterised what HierGIN can access;
the error was generalising from ``not in kNN-6 edge marginals'' to
``not in the data''.

\paragraph{Negative results.}
Patch-level scalar re-injection ($-$1.8\,pp to $-$3.0\,pp macro-F1) and
per-protein z-scoring of geometric features ($-$5.6\,pp) are both informative
failures.
The RI-GIN backbone gives a modest $+$0.2\,pp macro-F1 / $+$1.3\,pp balanced
accuracy over standard GIN on 8D features, consistent with a representation
regularisation effect rather than the large geometric gain seen in RIConv2.
LRF-GIN --- replacing on-the-fly RI scalars with precomputed 8D curvature-frame
edge features --- degrades macro-F1 by $-$3.4\,pp vs.\ RI-GIN.
Precomputed LRF frames from PCA are unstable on near-flat surfaces and cannot be
corrected at training time, confirming that on-the-fly invariant computation is preferable.
Three-level pooling (1\,000\,$\to$\,128\,$\to$\,32) using two independently-built
FPS assignments collapses to near-random performance (macro-F1\,=\,0.016);
nested assignment construction is required for multi-level pooling to be effective.

\paragraph{RI-HierNet validates the architecture direction.}
We implemented a multi-scale rotation-invariant hierarchical architecture
(RI-HierNet) that replaces HierGIN's kNN-6 node stage with four nested FPS +
ball-query set-abstraction layers using 6-D point-pair invariant features,
followed by a kNN-8 patch GIN over the resulting 64 centroids.
An earlier diagnostic prototype (GeoChem-MSRINet) produced unstable results
owing to curriculum and label-loading issues, but once resolved, the final
RI-HierNet run (\texttt{ri\_hiernet\_acc\_v4}) achieves
\textbf{accuracy\,=\,0.894 and weighted-F1\,=\,0.888} on test,
a $+$2.7\,pp accuracy improvement over HierGIN-14D (0.867).
This confirms the architecture hypothesis: multi-scale RI encoding of local
surface geometry yields a measurable gain over single-scale kNN graphs.
The remaining gap to RIConv2 RUN2 ($-$3.5\,pp accuracy) is attributable to
RI-HierNet's narrower encoder (512-D vs.\ RIConv2's 2048-D), shallower SA
stack (4 layers vs.\ 5), and absence of data augmentation.

\paragraph{Future work.}
RI-HierNet's architecture is validated but not yet exhausted.
Three targeted improvements could close the gap to RIConv2:
(1)~matching RIConv2's SA encoder width (2048-D output) and depth (5 layers);
(2)~applying the three augmentations used by RIConv2 (random point dropout,
random scale, random shift) which are straightforward to add to the data
pipeline;
(3)~removing label smoothing, which is absent in RIConv2's training recipe.
Longer term, gains beyond accuracy\,$\approx$\,0.93 will likely require
sequence features (e.g.\ ESM-2~\cite{esm2}) to resolve families whose
surfaces are structurally and chemically near-identical.
